\chapter{Stability and the Phase Plane}

Often we do not need the exact solution. We only need to know: \emph{does the system settle down, blow up, or oscillate?} This is \textbf{stability}. It is the main question in control, circuits and mechanics.

\section{Phase plane and critical points}

Consider an \textbf{autonomous} system (no explicit $t$ on the right side)
\[ x' = f(x,y), \qquad y' = g(x,y). \]
The $(x,y)$-plane is the \textbf{phase plane}. Each solution is a curve (\textbf{trajectory}) in it. A \textbf{critical point} (equilibrium) is a point where $f=g=0$, so the solution stays there forever.

\begin{defbox}{Types of stability}
\begin{itemize}
  \item \textbf{Stable:} trajectories that start near the critical point stay near it.
  \item \textbf{Asymptotically stable:} they stay near and also approach the point as $t\to\infty$.
  \item \textbf{Unstable:} some trajectories that start near it move away.
\end{itemize}
\end{defbox}

\section{Linear systems}

For $\mathbf{y}' = A\mathbf{y}$ with $\det A \neq 0$ the only critical point is the origin. Its type depends on the eigenvalues $\lambda_{1,2}$ of $A$ (Chapter 3).

\begin{center}
\renewcommand{\arraystretch}{1.35}
\begin{tabular}{@{}lll@{}}
\toprule
Eigenvalues & Type of origin & Stability \\
\midrule
real, both negative & stable node & asymptotically stable \\
real, both positive & unstable node & unstable \\
real, opposite signs & saddle point & unstable \\
complex, $\alpha<0$ & stable spiral & asymptotically stable \\
complex, $\alpha>0$ & unstable spiral & unstable \\
pure imaginary $\pm i\beta$ & center (closed loops) & stable, not asymptotically \\
\bottomrule
\end{tabular}
\end{center}

\begin{methodbox}{Quick test for a $2\times2$ matrix}
Let $p = \operatorname{trace}A$, $q = \det A$.
\begin{itemize}
  \item The origin is \textbf{asymptotically stable} exactly when $p<0$ and $q>0$.
  \item $q<0$: saddle (always unstable).
  \item $p^2 - 4q < 0$: complex eigenvalues (spiral if $p\ne0$, center if $p=0$).
\end{itemize}
This avoids computing the eigenvalues.
\end{methodbox}

\begin{derivbox}{Why the eigenvalues decide the type, and the trace--determinant test}
The general solution is $\mathbf{y} = \sum c_i\mathbf{v}_ie^{\lambda_it}$ (Chapter 3). For $\lambda = \alpha + i\beta$, $|e^{\lambda t}| = e^{\alpha t}$. So a term decays if $\operatorname{Re}\lambda<0$, grows if $\operatorname{Re}\lambda>0$, and keeps its size if $\operatorname{Re}\lambda=0$.
\begin{itemize}
  \item \textbf{Real, same sign:} both components move along the eigenvector directions, towards the origin (negative) or away from it (positive): a node.
  \item \textbf{Real, opposite signs:} one component decays and the other grows. Only the trajectories that start exactly on the decaying eigenvector reach the origin: a saddle.
  \item \textbf{Complex:} $e^{\alpha t}(\cos\beta t,\ \sin\beta t)$ describes a rotation with radius $e^{\alpha t}$: a spiral inward ($\alpha<0$), outward ($\alpha>0$), or a closed loop ($\alpha=0$, a center).
\end{itemize}
\textbf{Trace and determinant.} The characteristic equation of $A$ is $\lambda^2 - p\lambda + q = 0$ with $p = \lambda_1+\lambda_2 = \operatorname{trace}A$ and $q = \lambda_1\lambda_2 = \det A$. Then $\lambda = \dfrac{p\pm\sqrt{p^2-4q}}{2}$.
\begin{itemize}
  \item If $p^2<4q$: complex roots with $\operatorname{Re}\lambda = p/2$. Stable iff $p<0$.
  \item If $p^2\ge4q$ and $q<0$: the product is negative, so the roots are real with opposite signs: a saddle.
  \item If $p^2\ge4q$ and $q>0$: the roots are real with the same sign, and the sign is that of their sum $p$. Both negative iff $p<0$.
\end{itemize}
So both roots have negative real part exactly when $p<0$ and $q>0$.
\end{derivbox}

\begin{examplebox}{Classify three systems}
\begin{itemize}
  \item $A = \begin{pmatrix}1&2\\2&1\end{pmatrix}$: $\lambda = 3, -1$. Opposite signs, so a \textbf{saddle}, unstable. (Also $q = -3<0$.)
  \item $A = \begin{pmatrix}-1&2\\-2&-1\end{pmatrix}$: $\lambda = -1\pm2i$. A \textbf{stable spiral}. See the figure below.
  \item $my''+cy'+ky=0$ as a system. Using $u=y$, $v=y'$: $A = \begin{pmatrix}0&1\\-k/m&-c/m\end{pmatrix}$, so $p = -c/m$ and $q = k/m$. For $c>0$ the origin is asymptotically stable. It is a \textbf{node} if $c^2>4mk$ (overdamped), and a \textbf{spiral} if $c^2<4mk$ (underdamped). For $c=0$ it is a \textbf{center} (undamped oscillation). This is Chapter 2 in a new picture.
\end{itemize}
\end{examplebox}

\begin{center}
\begin{tikzpicture}
\begin{axis}[width=0.62\textwidth,height=6.2cm,trig format plots=rad,xlabel={$x$},ylabel={$y$},axis equal,
  xmin=-0.5,xmax=1.1,ymin=-0.7,ymax=0.5,grid=both,grid style={gray!20}]
\addplot[accent,thick,domain=0:5,samples=300]({exp(-x)*cos(2*x)},{-exp(-x)*sin(2*x)});
\addplot[only marks,mark=*,mark size=2pt,notec] coordinates {(1,0) (0,0)};
\end{axis}
\end{tikzpicture}

{\small Figure: trajectory of $x'=-x+2y$, $y'=-2x-y$ starting at $(1,0)$. It spirals into the origin.}
\end{center}

\section{Nonlinear systems: linearization}

\begin{methodbox}{Linearization at a critical point $(x_0,y_0)$}
\begin{enumerate}
  \item Solve $f=g=0$ to find the critical points.
  \item At each point compute the \textbf{Jacobian matrix}
  \[ J = \begin{pmatrix}f_x & f_y\\ g_x & g_y\end{pmatrix}_{(x_0,\,y_0)}. \]
  \item Find its eigenvalues (or use $p$ and $q$) and classify as in the table.
\end{enumerate}
\textbf{Rule (stated):} if the eigenvalues have non-zero real parts, the nonlinear system behaves like the linear one near the point. If the linear system is a center, the nonlinear system may be a center or a spiral, and we need more work.
\end{methodbox}

\begin{derivbox}{Linearization}
Let $(x_0,y_0)$ be a critical point and put $\xi = x - x_0$, $\eta = y - y_0$. Expand $f$ and $g$ in a Taylor series about the point, using $f(x_0,y_0) = g(x_0,y_0) = 0$:
\[ \xi' = f_x\,\xi + f_y\,\eta + (\text{terms of order } 2\text{ and higher}),\qquad \eta' = g_x\,\xi + g_y\,\eta + (\cdots). \]
Close to the point $\xi,\eta$ are small, so the quadratic terms are much smaller than the linear ones. We drop them and get the linear system $\begin{pmatrix}\xi\\\eta\end{pmatrix}' = J\begin{pmatrix}\xi\\\eta\end{pmatrix}$ with the Jacobian $J = \begin{pmatrix}f_x&f_y\\g_x&g_y\end{pmatrix}$ evaluated at the point.

\textbf{The pendulum equation.} A mass $m$ on a rod of length $L$ moves along a circle. The component of gravity along the path is $-mg\sin\theta$, the damping force is proportional to the speed, $-cL\theta'$, and the acceleration along the path is $L\theta''$. Newton's law gives $mL\theta'' = -mg\sin\theta - cL\theta'$, that is $\theta'' = -\dfrac gL\sin\theta - \dfrac cm\theta'$. With $\omega = \theta'$ and the values $g/L = 1$, $c/m = 0.5$ we get the system used below.
\end{derivbox}

\begin{examplebox}{Damped pendulum}
$\theta' = \omega$, $\omega' = -\sin\theta - 0.5\,\omega$. Critical points: $\omega=0$ and $\sin\theta=0$, so $\theta = 0$ (hanging down), $\theta=\pi$ (balanced up), and so on.

\textbf{At $(0,0)$:} $J = \begin{pmatrix}0&1\\-\cos0 & -0.5\end{pmatrix} = \begin{pmatrix}0&1\\-1&-0.5\end{pmatrix}$. $p=-0.5<0$, $q = 1>0$: asymptotically stable. $\lambda^2 + 0.5\lambda + 1=0$ gives $\lambda = -0.25\pm0.97i$, a \textbf{stable spiral}.

\textbf{At $(\pi,0)$:} $J = \begin{pmatrix}0&1\\-\cos\pi&-0.5\end{pmatrix} = \begin{pmatrix}0&1\\1&-0.5\end{pmatrix}$. $q = -1<0$, so a \textbf{saddle}: $\lambda = 0.78$ and $-1.28$. Unstable.

The pendulum hangs down stably and balances upright only unstably, as expected.
\end{examplebox}

\section{Stability of higher-order systems}

For a linear system or a transfer function, all solutions decay if and only if \textbf{every eigenvalue (pole) has a negative real part}. We can test this without finding the roots.

\begin{defbox}{Routh--Hurwitz test (up to degree 3)}
\begin{itemize}
  \item $s^2 + a_1s + a_0$: stable $\iff$ $a_1>0$ and $a_0>0$.
  \item $s^3 + a_2s^2 + a_1s + a_0$: stable $\iff$ $a_2,a_1,a_0>0$ \textbf{and} $a_2a_1 > a_0$.
\end{itemize}
\end{defbox}

\begin{derivbox}{Why the Routh--Hurwitz conditions are correct (degree 2 and 3)}
\textbf{Degree 2.} For $s^2 + a_1s + a_0$ the trace is $-a_1$ and the determinant is $a_0$ (the roots satisfy $s_1+s_2 = -a_1$, $s_1s_2 = a_0$). By the trace--determinant argument above, both roots have negative real part exactly when $a_1>0$ and $a_0>0$.

\textbf{Degree 3, necessity.} Suppose all roots have negative real part. A real cubic has a real root $-r$ with $r>0$, and the remaining quadratic factor $s^2 + bs + c$ has $b>0$, $c>0$ (degree 2 case). Expand $(s+r)(s^2+bs+c) = s^3 + (r+b)s^2 + (c+rb)s + rc$. So $a_2 = r+b$, $a_1 = c+rb$, $a_0 = rc$, all positive, and
\[ a_2a_1 - a_0 = (r+b)(c+rb) - rc = b\,(r^2 + rb + c)>0 . \]
\textbf{Degree 3, sufficiency.} Suppose $a_2,a_1,a_0>0$ and $a_2a_1>a_0$. A cubic with positive coefficients has no root $s\ge0$ (all terms are positive), so it has a real root $-r<0$. Factor $(s+r)(s^2+bs+c)$ with $b = a_2 - r$ and $c = a_0/r>0$. The same expansion gives $a_2a_1 - a_0 = b\,(r^2+rb+c)$, and $r^2 + rb + c = c + r(r+b) = c + ra_2>0$. Since $a_2a_1 - a_0>0$, we must have $b>0$. So the quadratic factor has $b,c>0$ and its roots have negative real part. Together with $-r<0$, all three roots lie in the left half-plane.
\end{derivbox}

\begin{examplebox}{Routh--Hurwitz}
$s^3 + 2s^2 + 3s + 1$: all positive and $2\cdot3 = 6>1$. Stable.
$s^3 + s^2 + 2s + 8$: all positive but $1\cdot2 = 2 < 8$. Unstable (has roots with positive real part).
\end{examplebox}

\begin{practicebox}
\begin{enumerate}
  \item Classify the origin for $x' = -2x$, $y' = -3y$.
  \item Classify the origin for $x' = y$, $y' = -4x$.
  \item Classify the origin for $x' = 2x + y$, $y' = x + 2y$.
  \item Use trace and determinant on $A = \begin{pmatrix}1&-3\\2&-4\end{pmatrix}$. Then find the eigenvalues.
  \item Is $s^3 + 2s^2 + s + 3$ stable?
  \item For $y'' + cy' + 4y = 0$, for which $c$ is the origin a stable spiral?
\end{enumerate}
\end{practicebox}

\begin{answerbox}
(1) $\lambda=-2,-3$: stable node. \quad (2) $\lambda = \pm2i$: center (stable, not asymptotically). \quad (3) $\lambda = 3, 1$: unstable node. \quad (4) $p=-3<0$ and $q = -4+6 = 2>0$, so asymptotically stable. $\lambda^2+3\lambda+2=0$ gives $\lambda=-1,-2$: stable node. \quad (5) $a_2a_1 = 2\cdot1 = 2 < 3$: unstable. \quad (6) The origin is a stable spiral if $c>0$ and $c^2 < 16$, so $0<c<4$.
\end{answerbox}
